\section{Reusable presentations for mathematical domains}

The purpose of a general framework is not merely to prove an abstract compiler
theorem.  It should tell an author how to replace familiar mathematical objects
by interfaces that the compiler can consume.  This section records a reusable
library of such presentations.

\subsection{Finite sequences, tables, and trees}

A finite sequence is coded by a length and an access function represented by a
single natural number.  Any standard primitive-recursive coding works; the
formalization should fix one globally.  The basic operations are:

\begin{itemize}
\item length, bounded lookup with a default value, and bounded replacement;
\item append, map by a fixed PR function, fold, filter by a Boolean predicate;
\item sorted insertion, duplicate removal, and finite membership;
\item tree child lists, parent pointers, and depth-first traversal.
\end{itemize}

Every theorem about such data is proved by a clocked fold with a quantifier-free
prefix invariant.  A useful design rule is to expose only the abstract laws in
node contracts and hide the chosen pairing code behind proved equations.  This
prevents application proofs from becoming arithmetic proofs about a particular
encoding.

\subsection{Integers and rationals}

Integers may be represented by pairs \((p,n)\) denoting \(p-n\).  Equality is
the natural-number equation
\[
        p+n'=p'+n.
\]
Addition and multiplication are polynomial operations on pairs.  The current
\path{IPRADerivations} library already implements this representation and
proves its ring laws in the quotient of terms by derivable equality.

Rationals may be represented by signed numerators and positive denominators.
There are two useful equality notions.

\begin{enumerate}
\item \emph{Extensional equality}: \((a,b)\sim(c,d)\) iff \(ad=cb\), with
positivity guards for \(b,d\).  This avoids gcd normalization and is often best
for proofs.
\item \emph{Canonical equality}: require positive denominator and gcd one.
This gives decidable literal equality but needs a verified gcd and
normalization development.
\end{enumerate}

For a general framework, use extensional equality in local algebra and invoke a
canonicalizer only when an external format requires it.  Division by a rational
is guarded by a nonzero numerator certificate.  Order comparisons are reduced
to integer inequalities after multiplying by positive denominators.

\subsection{Dyadic error arithmetic}

Approximation proofs frequently need only dyadics \(m/2^q\).  A dyadic is a
signed integer code and a nonnegative exponent.  Addition aligns exponents;
multiplication adds exponents.  The inequality
\(|x|\le2^{-q}\) is decided after shifting integers.

A regular name can therefore be stated entirely with integer equations and
inequalities:
\[
   |2^{q'}a(q)-2^q a(q')|\le 2^{q+q'-\min(q,q')}\,\times\,\text{the chosen tolerance constant}.
\]
In practice it is clearer to define a fixed dyadic distance library and prove
its algebra once.

\subsection{Represented real numbers}

\begin{definition}[Regular rational name]
A regular name is a PR function \(a:\N\to\Q\) satisfying
\[
       |a(q)-a(r)|\le 2^{-q}+2^{-r}
       \qquad\text{for all }q,r.
\]
A faster convention may use \(|a(q)-a(r)|\le2^{-q}\) when \(r\ge q\).
\end{definition}

The object-language presentation uses only rational codes and the displayed
finite inequalities.  The external adequacy theorem constructs a Cauchy real.
Operations are contracts with explicit demands.

\paragraph{Addition.}
To output precision \(q\), request both inputs at \(q+2\), add the rationals,
and allocate the two input errors plus exact rational arithmetic error.

\paragraph{Multiplication.}
Multiplication needs magnitude bounds.  A contract includes integers
\(M_x,M_y\) with \(|x|\le M_x\), \(|y|\le M_y\).  Choose input precision
large enough that
\[
    M_x\varepsilon_y+M_y\varepsilon_x+
    \varepsilon_x\varepsilon_y\le2^{-q}.
\]
The inequality is checked rationally.

\paragraph{Reciprocal.}
Require an apartness certificate \(|x|\ge2^{-g}\).  Request enough precision to
preserve a lower bound on the rational approximation and to meet the standard
error estimate for inversion.  The operation is totalized with a default when
the guard fails; correctness is conditional on the guard.

\paragraph{Limits.}
A sequence of names has a limit contract only when supplied with a Cauchy
modulus or a metastable bound adequate for the desired conclusion.  Merely
asserting that the sequence converges classically is insufficient for a
uniform producer.

\subsection{Complex numbers}

Represent a complex number by a pair of represented reals, or by Gaussian
rational approximations.  Addition and multiplication are polynomial in the
components.  Inversion is guarded by a lower bound on
\(x^2+y^2\).  Absolute value can be represented by a squared norm when square
roots are unnecessary; this often simplifies polynomial arguments.

Root isolation, argument-principle counts, and contour estimates should use
rational rectangles or disks with certified separation from zero.  The branch
of a square root or logarithm is a relational object unless a domain and cut
certificate select it.

\subsection{Polynomials and formal series}

A polynomial over a coded coefficient domain is a finite coefficient list.
Degree is bounded by list length; evaluation is a Horner trace.  Addition,
multiplication, formal derivative, coefficient truncation, and composition to
a fixed degree are primitive recursive.

A formal power series is not stored as an infinite list.  It is a coefficient
producer \(a(n)\) together with finite contracts.  A query \(N\) requests the
prefix \((a(0),\ldots,a(N))\).  Cauchy products and formal composition use
bounded double sums.  Analytic evaluation additionally needs a radius and a
remainder bound; formal identities do not.

\begin{proposition}[Finite formal-series algebra]
Every identity between coefficients up to a requested order \(N\), built from a
fixed finite expression using addition, multiplication, derivative, integral
with explicit denominator guards, and composition with zero constant term, is
computable by a bounded PR trace and reducible to polynomial coefficient
identities.
\end{proposition}

\begin{proof}
At coefficient index \(n\le N\), each operation uses only finitely many indices
at most \(n\).  Evaluate the expression tree by a topological graph and the
coefficient recurrences by bounded sums.  Correctness follows by induction on
the expression and on \(n\).  Denominator divisions use the explicit fact that
\(n+1>0\).
\end{proof}

\subsection{Matrices and linear algebra}

For a fixed dimension \(d\), a matrix is a table of \(d^2\) entries.  Matrix
addition and multiplication are finite sums.  Determinants may be defined by a
Leibniz sum over the finite permutation list, a Laplace recursion with explicit
fuel, or a division-free elimination algorithm.  The best representation
depends on the theorem.

For polynomial identities such as Cayley--Hamilton at fixed \(d\), a
permutation or adjugate proof is often easiest to compile: all sums are bounded
by constants depending on \(d\), and the remaining algebra is polynomial.
For variable \(d\), the theorem remains plausible in iPRA, but the formal
library needs generic finite permutations, signs, matrices, determinants, and
polynomial matrices.  The framework supplies the control and proof-composition
pattern, not those domain lemmas for free.

Gaussian elimination is guarded by pivot nonzeroness.  To avoid branch
instability, one may enumerate all pivot choices, retain the successful
relations, and observe rank, determinant, or solution-space dimension rather
than a particular row-echelon form.  Exact rational elimination is appropriate
when all entries are exact; approximate matrices require separation bounds.

\subsection{Finite groups, rings, and modules}

A finite algebraic structure is a carrier bound, operation tables, and Boolean
axiom checks.  A homomorphism is a finite table satisfying bounded equations.
Subgroups, ideals, kernels, and images are bit vectors with closure checks.

For finitely generated infinite structures, use finite presentations and query-
dependent normal forms only when an explicit terminating rewriting or bounded
word problem is available.  A group presentation alone does not supply a PR
normal-form algorithm.

Modules over exact rings can be represented by finite matrices.  Smith normal
form is a descent algorithm with explicit coefficient-growth and measure bounds
needed for a uniform iPRA proof.  For a fixed concrete matrix, a proposed Smith
form plus unimodular multipliers is a finite certificate easily verified by
matrix multiplication and determinant checks.

\subsection{Simplicial and cellular data}

A finite simplicial complex is a finite list of vertex subsets closed under
faces, or a Boolean face predicate on a bounded vertex set.  Chain groups are
finite free abelian groups, boundary maps are integer matrices, and the
identity \(\partial^2=0\) is a finite signed cancellation theorem.  Homology
invariants can be computed from Smith certificates.

A triangulated or cubical approximation to a space is only a finite model.  A
theorem relating it to an ideal topological space needs explicit approximation
or equivalence data: a simplicial approximation, deformation retraction,
nerve-cover certificate, or mesh/regularity theorem.  The framework cleanly
separates the finite homological computation from this producer.

\subsection{Finite metric and topological observations}

A finite metric observation contains a finite point list and rational upper and
lower bounds for distances.  Triangle inequalities and separation properties
are bounded checks.  Nets, covers, Rips complexes, and Cech complexes at a
rational scale are finite constructions.

For a compact represented space, a query may request a finite
\(2^{-q}\)-net plus a covering certificate.  A continuous map is represented by
an observation transformer and a modulus.  Uniform continuity is therefore not
silently obtained from compactness; the modulus or finite cover argument is
part of the contract.

\subsection{Finite categories and diagrams}

A finite category is a finite object set, finite morphism set, source and target
tables, identity table, partial composition table, and bounded proofs of the
category laws.  A finite functor and natural transformation are finite tables
with bounded naturality checks.  Commutative diagrams become equations in
finite hom-sets.

For large or infinite categories, a theorem normally uses only a finite
diagram.  The presentation should therefore encode the diagram and the local
universal-property witnesses actually needed.  ``All objects'' statements are
ordinary quantified formulas, but a construction that searches over all
objects is not available without a bounded presentation.

Yoneda-style arguments can be expressed extensionally: a natural
transformation is determined by its component at the representing object, and
the proof is finite symbolic composition.  Choosing representatives of
isomorphism classes is unnecessary.

\subsection{Probability distributions}

A finite rational distribution is a finite list of nonnegative rational weights
summing to one.  Expectation of a bounded rational variable is a finite sum.
Convolution, moments to a fixed order, and finite Markov chains are PR table
algorithms.

An infinite law is represented by finite cylinder probabilities or rational
approximations with consistency and error bounds.  Limit theorems require
explicit finite error estimates.  A qualitative weak-convergence statement
without a modulus may be reconstructed metastably, but should not be advertised
as a uniform rate.

\subsection{A general recipe}

For a new domain, the recommended order is:

\begin{enumerate}
\item choose the smallest finite observations used by the proof;
\item define validity and restriction before defining operations;
\item prove coding and bounded-search lemmas once;
\item classify each operation as exact, guarded, relational, or approximate;
\item attach backward demand and resource effects;
\item prove local contracts in \(\Prf\;\PRAAx\) or \(\KPrf\);
\item use the graph compiler only after the local library is stable;
\item state an external adequacy theorem connecting observations to the
classical object.
\end{enumerate}

This order keeps representational decisions from being buried inside the proof
of a headline theorem.
